\section{The original one-character trace and length argument}
\label{sec:one-character-v48}

This chapter reconstructs the implication whose source name is
\texttt{oneChar}: the exact \texttt{InpZ} together with
\texttt{AutomorphyUnramifiedCoefficients} implies \texttt{OneChar}.
The latter automorphy assertion remains an input.  It supplies one
marked element and three properties of the same deformation problem;
it is not replaced by three unrelated existence assertions.
All source references are at
\texttt{7cbf7a0c32fca1badf05931879a37c3f8a2b03b4}.
The original local period-ring, finite Galois-cohomology, class-field,
and Poitou--Tate foundations are those already recorded in BX1--BX8,
DX1--DX3 and KB8.  No general crystalline deformation-ring theorem is
added as a hypothesis below.

Write \(G=G_{\mathbf Q}\), \(H=G_K\), and
\(\epsilon=\epsilon_\ell\).  The coefficient datum is the original
\(c\), with discrete valuation ring \(\mathcal O\), fraction field
\(E\), and specified inclusion into
\(\overline{\mathbf Q}_\ell\).  The datum \(GD\) includes the
original imaginary quadratic field, \(e\ge1\), odd \(\ell\),
and all its conductor, parity, order, coprimality and residual
inertia conditions.  We preserve that predicate throughout.
Put \(\Psi=\epsilon^e\psi\), let \(\eta\) be the quadratic
character of \(K/\mathbf Q\), and set
\[
 w=(1,0,-1),\qquad \operatorname{wt}(i,j)=w_i-w_j,
 \qquad T(t)=\operatorname{diag}(t,1,t^{-1}).
 \tag{OCH.1}
\]
The disconnected group is
\(\mathcal G=\mathrm{GL}_3\rtimes\{\pm1\}\), where the
nonidentity component acts by
\(\theta(A)=J A^{-t}J\), with \(J\) the reversing matrix.
The base representation has matrix \(T(\Psi(g))\) and component
\(\eta(g)\).  Its adjoint action in that component is
\(-JX^tJ\), not ordinary conjugation alone.
All matrix and coefficient maps below are actual maps, before passage
to isomorphism classes.

\subsection{Positive local cohomology and the specified ordinary component}

% BEGIN OCH1
\paragraph{OCH1. Exactness for positive finite-coefficient representations.}
\label{claim:OCH1}
Let \(L/\mathbf Q_\ell\) and \(E/\mathbf Q_\ell\) be finite.
Consider a de Rham \(E\)-representation with a filtration whose
successive characters, on an open inertia subgroup, are positive
integral powers of \(\epsilon\).  On this category the crystalline
finite-cocycle functor has the following properties: it maps a
surjection onto the finite cocycles of the quotient; and a cocycle
in a subrepresentation is finite if its image is finite.  The second
assertion assumes positivity of both representations.  These
statements allow nilpotent coefficient algebras, viewed as finite
\(E\)-vector spaces.

\emph{Proof.}
Use BX1's actual period rings and BX5's fundamental exact sequence.
For a positive character, pass to a finite extension killing its
finite inertia part.  It is \(\epsilon^a\) times an unramified
character, \(a\ge1\).  The unramified value has absolute value one:
the character is continuous on a compact group, so its image and its
inverse image are bounded.  The eigenvalues of crystalline Frobenius
therefore have nonzero slope \(-a\).  In particular \(1-\varphi\)
is invertible on its crystalline periods.  The same holds for a
successive extension, by an upper triangular matrix calculation;
finite restriction does not introduce a fixed vector.  Positive
weights also give \(\mathrm{Fil}^0D_{\rm dR}=0\).
The fundamental exact sequence consequently identifies
\(H^1_f(L,V)\) with \(D_{\rm dR}(V)\).  De Rham periods are
exact on de Rham representations, as follows by choosing a period
basis and tensoring the short exact sequence with \(B_{\rm dR}\).
Thus the asserted surjectivity and injectivity hold on classes.
For surjectivity on cocycles, lift the desired class, compare its
image with the given cocycle, and lift the vector producing their
coboundary difference.  For injectivity no adjustment is needed.
Scalar restriction, tensor coefficient retractions, and restriction
of periods preserve the original local kernels by KTR1.  This proves
all statements for the supplied coefficient algebras, not merely for
semisimple characters.
% END OCH1

% BEGIN OCH2
\paragraph{OCH2. The only positive-character obstruction is the cyclotomic annihilator.}
\label{claim:OCH2}
Let \(B\) be a finite local \(E\)-algebra with residue \(E\),
and let \(\kappa:G_L\to B^\times\) restrict to \(\epsilon^a\)
on inertia, \(a\ge1\).  Choose \(F\) generating the quotient
by inertia.  Every continuous cocycle can be written
\[
                  z=z_f+z_s,
 \qquad (\kappa(F)-\epsilon(F))z_s(g)=0,
 \tag{OCH.2}
\]
where \(z_f\) is finite.  If \(a\ge2\), one may take
\(z_s=0\).

\emph{Proof.}
Set \(d=\kappa(F)-\epsilon(F)\) if \(a=1\), and \(d=1\)
otherwise.  The ideal generated by
\(\kappa(g)-\epsilon(g)\), for all \(g\), is \((d)\).
For \(a=1\) divide the two characters and use the dense powers of
\(F\); the ideal \((d)\) is closed because \(B\) is finite
 dimensional.  For \(a>1\), inertia contains an element whose
\(\epsilon^{a-1}\)-value differs from one; the resulting nonzero
scalar in \(E\) is a unit.
Let \(W=\operatorname{ann}_B(d)\).  Both \(W\) and \(B\)
with the \(\kappa\)-action are positive representations.  The
local Euler characteristic and Tate duality in BX8 give
\[
 \dim_E Z^1(V)-\dim_E Z^1_f(V)
       =\dim_E H^0(V^*(1))
 \tag{OCH.3}
\]
for either of them; the finite-cocycle dimension here is the
Bloch--Kato dimension formula, with the same actual coboundaries
subtracted on both sides.  The last invariant space for \(B\)
is the dual of \(B/dB\), of dimension \(\dim_E W\).
On \(W\) the action is cyclotomic, so its corresponding dimension
is also \(\dim_E W\).  OCH1 makes the map
\(Z^1(W)/Z^1_f(W)\to Z^1(B)/Z^1_f(B)\) injective;
(OCH.3) makes it surjective.  Subtract a cocycle valued in \(W\)
representing the given quotient class.  This is (OCH.2).
When \(d=1\), \(W=0\).
% END OCH2

% BEGIN OCH3
\paragraph{OCH3. Crystalline triangular lifts over every finite Artin quotient.}
\label{claim:OCH3}
For the original local datum \(D:\texttt{Dat3}\), every upper
triangular crystalline deformation of the base representation lifts
through every surjection of finite local \(E\)-algebras inducing
identity on the residue field.  The component character remains
\(\eta\); its inertia diagonal remains
\((\Psi,1,\Psi^{-1})\).

\emph{Proof.}
Write an upper triangular matrix as a diagonal torus matrix times
\[
 U(x,y,z)=
 \begin{pmatrix}1&x&z+xy/2\\0&1&y\\0&0&1\end{pmatrix}.
 \tag{OCH.4}
\]
The diagonal entries constitute the source's torus cocycle.
After division by their fixed inertia characters, their infinitesimal
changes are unramified.  Indeed each successive square-zero diagonal
variation is a finite cocycle in a weight-zero character module.
The weight-zero Bloch--Kato formula identifies its finite cohomology
with the unramified classes.  The finite local component action is
retained: apply that formula after the finite extension trivializing
it, then take invariants by averaging over the finite quotient.
The same argument respects the disconnected involution on the torus.
Induction through the nilpotent filtration therefore fixes the entire
inertia character, not only its residual value.
Lift their logarithms through the surjection: logarithm and exponential are finite polynomials on the
nilpotent ideals.  On the unramified quotient the prescribed value
at Frobenius determines the character continuously; the involution
on the diagonal is respected by taking its \(+\) and \(-\)
parts, which is permissible since \(2\) is invertible.
This lifts the actual torus, including its disconnected action.

The pair \((x,y)\) is a cocycle in the two-dimensional
superdiagonal module.  Its residual characters are \(\Psi\)
and \(\eta\Psi\); the filtration by powers of the coefficient
maximal ideal makes this module positive.  Crystallinity of the
original triangular representation says exactly that this pair is
finite.  OCH1 lifts that finite cocycle, and it can be adjusted by a
lifted coboundary to have precisely the prescribed image.
The multiplication law of (OCH.4) now prescribes the coboundary of
\(z\).  Its obstruction is in \(H^2\) of the central character
\(\eta\kappa_0/\kappa_2\), whose inertia weight is \(2e\ge2\).
Local duality makes this \(H^2\) zero: a dual invariant would have
inertia weight \(1-2e\ne0\).  Thus an actual continuous \(z\)
exists.  This assertion follows equally from continuous cochains:
a finite-dimensional cohomology class that is zero has a continuous
cochain primitive.

The resulting representation is crystalline.  Over the open subgroup
\(D.HL=\ker\eta\cap\ker\tau\), its first two columns give the
finite extension determined by \(x\).  Twisting that extension by
\(\kappa_2^{-1}\), its last column is an extension with quotient
character of weight \(e\) and kernel character of weight \(2e\).
The kernel has \(H^1=H^1_f\) and \(H^2=0\), by OCH2 and
local duality.  Consequently the last-column cocycle is finite
exactly when its \(y\)-coordinate is finite.  This holds by
construction.  Finite-index descent uses the average of a period
primitive over the finite quotient; the characteristic-zero
coefficient field allows division by its order.
Finally, the difference between the constructed \(z\) and the
specified \(z\) after reduction is a finite central cocycle.
Apply OCH1 once more to lift that difference and add it.  Every
matrix entry now reduces to the original one.  This proves the
literal \texttt{triLift}, including its all-surjections quantifier.
% END OCH3

% BEGIN OCH4
\paragraph{OCH4. Every crystalline deformation of the base is ordinary.}
\label{claim:OCH4}
Over a finite local \(E\)-algebra \(B\) with residue \(E\),
a continuous deformation of the specified base representation is
crystalline on \(D.HL\) if and only if it is a conjugate of an
upper triangular crystalline deformation by a lower unipotent matrix
congruent to one.  The crystalline deformation functor is formally
smooth on these finite Artin algebras.

\emph{Proof.}
For dual numbers write the deformation as
\((1+\varepsilon\gamma(g))r_0(g)\).  It is crystalline exactly
when \(\gamma\) is a finite adjoint cocycle: apply the period
criterion to the extension of the base representation by itself.
The negative matrix-weight parts have inertia weights \(-e\)
and \(-2e\).  Their \(H^1_f\) vanishes by the character
Bloch--Kato formula in BX5.  Hence these parts are coboundaries.
Conjugating by \(1+\varepsilon X\), with \(X\) strictly lower
triangular, kills precisely them.  The diagonal and upper entries
remain unrestricted finite cocycles.  This proves the assertion
for dual numbers.

For general \(B\), use induction on its length.  A nonzero last
power of its maximal ideal contains a line \(E\varepsilon\)
annihilated by that ideal.  Apply induction modulo this line.
OCH3 lifts the triangular representative; the three entries of its
lower conjugator lift freely.  The difference between this lift and
the supplied representation is a dual-number deformation.  Its
crystalline condition really does follow from the given two
representations: their amalgam lies in
\(B\times_{B/(\varepsilon)}B\), and its underlying representation
is a subrepresentation of their direct sum.  Crystalline
representations are closed under subrepresentations.  Identifying
this fibre product with \(B\times_E E[\varepsilon]/(\varepsilon^2)\)
turns the difference into exactly the preceding finite cocycle.
Eliminate its lower part and glue its upper part.  This completes
the induction.  Conversely conjugation preserves periods, and
ordinary representatives were required to be crystalline.
Lift an ordinary representative by OCH3 and its lower conjugator
entry by entry to prove formal smoothness.  This is the full
\texttt{ordEq}/\texttt{crLift} argument, including nilpotent
coefficient directions.
% END OCH4

% BEGIN OCH5
\paragraph{OCH5. The saturated ordinary ideal: all its finite-algebra points are crystalline.}
\label{claim:OCH5}
Let \(U\) be the actual local universal ring and let \(I_Y\)
be its original saturated ordinary ideal.  Every map from
\(U\) to a finite \(\mathbf Q_\ell\)-algebra annihilating
\(I_Y\) gives a representation crystalline on \(D.HL\).

\emph{Proof.}
Here is the precise ideal, so that no unnamed crystalline component
is substituted.  Choose the original inertia element \(s_0\in D.HL\)
with infinite cyclotomic value and the original Frobenius \(F\)
of \(D.HL\).  Put \(\psi_i=\Psi^{w_i}\) and
\[
 c_i=\prod_{j\ne i}(\psi_i(s_0)-\psi_j(s_0)),\qquad
 R_i(s)=r(s)-\psi_i(s)I.
\]
The ordinary ideal \(I_{\rm ord}\) imposes the characteristic
polynomials \(\prod_i(X-\psi_i(s))\) on inertia and every ordered
product \(R_0(s)R_1(s')R_2(s'')=0\).
Set \(t_i=\operatorname{tr}(R_{i+1}(s_0)R_{i+2}(s_0)r(F))\),
where the indices are cyclic, and
\[
 d_i=c_{i+1}t_i-\epsilon(F)c_i t_{i+1}\quad(i=0,1),
 \qquad \Delta=\ell d_0d_1,
 \qquad I_Y=I_{\rm ord}:\Delta^\infty.
 \tag{OCH.5}
\]
These are the source's \texttt{cc,tr,dd,del,IY}.

First reduce to a local finite target \(B\).  In
\((U/I_Y)[1/\ell]\), multiplication by \(\Delta\) is
injective, by saturation.  If its map to \(B\) kills a power
\(\mathfrak m^a\) of the inverse image of the maximal ideal,
Artin--Rees for the submodule \(\Delta A\subset A\) gives
\(b\ge a\) such that
\(\Delta z\in\mathfrak m^b\) implies
\(z\in\mathfrak m^a\).  The quotient
\(B'=A/\mathfrak m^b\) is a finite local
\(\mathbf Q_\ell\)-algebra mapping to \(B\), and that map
kills \(\operatorname{ann}_{B'}(\Delta)\).  Its finiteness
follows from the finite residue field extension and the finite
successive quotients of the maximal-ideal filtration.  Thus it is
an actual finite thickening, even when \(\Delta\) vanishes
at the given point.

Over \(B'\), the numbers \(c_i\) are nonzero coefficient-field
scalars, hence units.  The spectral projectors
\(c_i^{-1}R_{i+1}(s_0)R_{i+2}(s_0)\) have rank one.
The ordered-product equations make their partial sums the successive
stable flag: the rightmost factor kills the last quotient, then the
middle factor the next, then the leftmost factor the first.
Multiplying the equations by the projectors proves this for each
inertia element.  This ordered inertia flag is canonical.  Its first
line is the common \(\psi_0\)-eigenline: the projector for
\(s_0\) makes any such line equal to the displayed rank-one
summand.  In the quotient the same argument singles out the
\(\psi_1\)-line, since the three values at \(s_0\) are
distinct.  The characters \(\psi_i\) extend to the
local group, so conjugation carries each of these labelled inertia
submodules to one with the identical label.  Normality of inertia
and uniqueness then preserve both steps of the flag.
In a flag basis the representation on \(D.HL\) is upper triangular, with characters
\(\kappa_i\) whose inertia restrictions are \(\psi_i\).
Its adjacent normalized entries are cocycles of
\(\kappa_0/\kappa_1\) and \(\kappa_1/\kappa_2\).
Moreover \(t_i=c_i\kappa_i(F)\), so \(d_i\) is a unit
multiple of \(\kappa_i(F)/\kappa_{i+1}(F)-\epsilon(F)\).
OCH2 says that the nonfinite part of each adjacent cocycle takes
values in the annihilator of this last element.  The map to \(B\)
kills those annihilators, since it kills the annihilator of their
product \(\Delta\).  Both adjacent cocycles in \(B\) are
therefore finite.  The last-column argument of OCH3 makes the full
rank-three representation crystalline.  Conjugation does not affect
this conclusion.  A general finite algebra is a finite product of
local ones; periods and their spanning criterion respect this
product.  This proves the stated all-target assertion.
% END OCH5

% BEGIN OCH6
\paragraph{OCH6. Crystalline Artin deformations annihilate the same saturated ideal.}
\label{claim:OCH6}
Every finite Artin deformation reducing to the prescribed base and
crystalline on \(D.HL\) annihilates \(I_Y\), including points
where \(\Delta=0\).

\emph{Proof.}
OCH4 gives an ordinary flag.  Its inertia characters are fixed:
a crystalline deformation of a fixed character has unramified
infinitesimal variation, by the weight-zero Bloch--Kato formula.
Thus the characteristic-polynomial and ordered-product equations
vanish, and the map annihilates \(I_{\rm ord}\).
To pass to the saturation, suppose \(\Delta^k a\in I_{\rm ord}\).
Let \(B\) be the finite Artin target, with maximal ideal
\(\mathfrak n\) satisfying \(\mathfrak n^h=0\).
Choose \(M=2kh+1\),
\(C=E[t]/(t^M)\), and \(T=B[t]/(t^M)\).
Evaluation at zero and reduction modulo \(\mathfrak n\) give
\(p_1:T\to B\) and \(p_2:T\to C\); together they surject onto
\(B\times_E C\).
Twist the base torus by the unramified character taking the original
arithmetic Frobenius to \(1+t\).  It is crystalline and reduces
to the base.  Substitution in (OCH.5) gives
\[
                   \Delta=t^j u,\qquad 0\le j\le2,
                       \quad u\in C^\times.
 \tag{OCH.6}
\]
Indeed each adjacent difference is either a unit at \(t=0\)
or has a nonzero linear coefficient; the positive power of
Frobenius used in the source multiplies that coefficient by a
nonzero integer in \(E\).  This also explains why the potentially
critical case \(e=1\) is included.

Amalgamate this twist with the supplied deformation over
\(B\times_E C\).  The amalgam is crystalline as a
subrepresentation of their direct sum.  Lift it to \(T\) by
OCH4, and use the universal property of the actual local ring.
It annihilates \(I_{\rm ord}\), so its images \(P=\Delta^k\)
and \(y=a\) satisfy \(Py=0\); also
\(p_2(P)=t^{jk}u^k\).  Lift the unit \(u^k\) to a unit
\(V\in T\).  Then \(P=t^{jk}V+Q\), with
\(Q\in\mathfrak nT\).  Consequently
\(t^{jk}y=-QV^{-1}y\), and iteration gives
\(t^{jkh}y=0\).  Since \(jkh<M\), the coefficient of
\(t^{jkh}\) in this equality is the constant coefficient of
\(y\).  It vanishes: \(p_1(y)=0\), the required original
image of \(a\).
For a residue field larger than \(E=\operatorname{Frac}\mathcal O\),
first take the subalgebra of elements whose residue lies in \(E\).
The original universal map lands there because its reduction is the
fixed base map.  Apply the preceding argument there and include it
back into the original target.  No common coefficient enlargement
or noncriticality assumption has been imposed.
% END OCH6

% BEGIN OCH7
\paragraph{OCH7. The literal local component has a twelve-variable weighted chart.}
\label{claim:OCH7}
The source's ideal \(I_t\), the intersection of the minimal primes
of \(\sqrt{I_Y}\) through the base, is a single prime.  Its
characteristic-zero completion at the base is a formal power-series
ring in twelve variables.  The variables can be given weights with
multiplicities
\[
 \begin{array}{c|rrrrr}j&-2&-1&0&1&2\\\hline
                   \#X_j&1&2&3&4&2\end{array}.
 \tag{OCH.7}
\]
There is an actual map
\(U\to E[[X_1,\ldots,X_{12}]]\) with kernel \(I_t\),
constant coefficient the specified base map, and every universal
matrix entry homogeneous of its original matrix weight.  It has
both the linear-coordinate surjectivity and the denominator-cleared
cotangent injectivity required by \texttt{LocChart}.

\emph{Proof.}
OCH5--OCH6 identify the completed local functor of \(U/I_Y\)
with the crystalline deformation functor, on every finite Artin
algebra.  OCH4 proves formal smoothness.  Choose a basis of its
cotangent space and lift it successively modulo powers of the maximal
ideal.  Formal smoothness lifts the resulting maps at each stage;
the inverse limit is a surjection from a power-series ring.  A
relation of least degree would prevent lifting through the quotient
by that relation and the next maximal-ideal power.  Thus the kernel
is zero.  This proves the power-series assertion directly from the
lifting property.  Completion is faithfully flat on this Noetherian
local ring, so the local ring is a domain; there is exactly one
minimal prime through the base.  Pulling it back to \(U\) gives
\(I_t\), the kernel of the displayed map.  Equivalently one can
clear a single localization denominator \(y\), with nonzero base
value, to obtain \(yI_t\subset I_Y\).

The tangent space is the space of finite adjoint cocycles at \(r_0\).
Its weight \(j\) has matrix dimension \(1,2,3,2,1\).
The framed Bloch--Kato dimension formula is
\(\dim Z^1_f(V)=\dim V+\dim D_{\rm dR}(V)/\mathrm{Fil}^0\):
subtract \(H^0\) when forming coboundaries and add the same
\(H^0\) in the finite-cohomology formula.  The extra term is
zero for \(j\le0\) and is respectively \(2,1\) for
\(j=1,2\).  This gives (OCH.7), including the case where
\(\eta\) is nontrivial unramified locally.

Universal entry deviations span the cotangent space, since the
universal ring is constructed from those entries.  Select twelve
homogeneous deviations forming an \(E\)-basis.  The formal inverse
function argument changes the coordinates to precisely their images.
Conjugation by \(T(u)\) acts on each selected coordinate by
\(u^j\).  Choose an infinite-order \(u\in\mathcal O^\times\).
The equality of formal series under this action forces every monomial
of a matrix entry to have its asserted weight: unequal powers of
\(u\) are distinct.  Thus the map is graded.
If an element has zero constant and linear coefficients, its image
is in the square of the formal maximal ideal.  Faithful flatness,
then localization, puts a multiple by an element \(y\) with
nonzero base value in \((\ker x)^2+I_t\).  Modulo this square,
\(y=x(y)+(y-x(y))\); multiplying once more replaces \(y\)
by a nonzero scalar from \(\mathcal O\).  This proves the
literal denominator-cleared injectivity.  Clearing denominators in
the chosen twelve linear coordinates proves their required
surjectivity.  All assertions concern the original integral prime
\(I_t\), not an ideal chosen only after extending coefficients.
% END OCH7

% BEGIN OCH8
\paragraph{OCH8. Conjugation stability over every original complete local target.}
\label{claim:OCH8}
The ideal \(I_t\) is stable under every near-identity conjugation
in the exact \texttt{ConjStable} sense.  Together with OCH7 this
proves the full original \texttt{LocGood} assertion.

\emph{Proof.}
Let \(R=U/I_t\), a domain, and form
\(R[[Y_{ij}]]\).  The matrix \(1+Y\) is invertible and
congruent to one.  Conjugating the universal representation produces
an actual universal map
\(\alpha:U\to R[[Y]]\).  The ordinary equations and the
trace expressions defining \(\Delta\) are invariant under
conjugation.  Moreover \(\Delta\ne0\) in \(R\): it is a
nonzero divisor on \(U/I_Y\), so belongs to no minimal prime
of that ring.  Hence \(\alpha\) kills \(\sqrt{I_Y}\).
Its kernel is prime because its target is a domain.  Evaluation at
\(Y=0\) is the original quotient map, by uniqueness in the
universal property; consequently \(\ker\alpha\subset I_t\).
Minimality of \(I_t\) over \(\sqrt{I_Y}\) forces equality.
Now, for any original complete local target, substitute the entries
of \(h-1\) for the \(Y_{ij}\), and the supplied map for
\(R\).  Adic convergence gives the map on the full power-series
ring.  Its composite with \(\alpha\) is the actual conjugated
representation and annihilates \(I_t\).  This proves the universal
all-target statement, including targets with nilpotents or torsion.
% END OCH8

\subsection{The actual global ring and its graded model}

% BEGIN OCH9
\paragraph{OCH9. The datum, the local universal ring, and the normalized global ring exist.}
\label{claim:OCH9}
Every original \(GD\) datum admits the exact \texttt{Dat} and
\texttt{ULoc}.  For each supplied marked \(g_1\), the original
normalized ring \(S:\texttt{Dat.Ring}\) exists and represents
actual continuous lifts with the prescribed multiplier, local
\(I_t\)-condition, killed inertia, and normalization.

\emph{Proof.}
The quadratic extension gives \(\eta\) with kernel \(G_K\).
Choose an embedding of the algebraic closure into \(\mathbf C\)
to specify \(c\), and the original embedding into
\(\overline{\mathbf Q}_\ell\) to specify the local inclusion
\(\iota_\ell:G_{\mathbf Q_\ell}\to G\).  The equivariant
extension of embeddings identifies its range with the chosen
 decomposition group and its inertia with the chosen inertia.
Since \(\ell\nmid\operatorname{disc}K\), \(\eta\) is
unramified locally.  Put \(\tau=\operatorname{comp}_\ell\psi\)
on local inertia.  Its finite image gives an open kernel, and
\(\Psi|_{I_\ell}=\epsilon^e\tau\).  Parity in \(GD\)
gives \(\Psi(c)=1\).  These choices satisfy the actual
\texttt{Dat3} and \texttt{Dat} fields; none changes the supplied
coefficient datum or character.

Define \(\mathrm{Kill}\) to be the union, at all places away
from \(\ell\), of the kernels of the residual representation
inside inertia.  The prime-to-\(\ell\) order of \(\psi\)
implies that the base representation itself is one on this set:
reduction is injective on its prime-to-\(\ell\) roots of unity,
and the cyclotomic character is unramified there.  Choose a finite
set of places containing \(\ell\), the conductor and the
quadratic discriminant.  The quotient by the closed normal closure
of \(\mathrm{Kill}\) is a quotient of \(G_{\mathbf Q,S}\).
The finite-cohomology hypothesis gives its \(\ell\)-finiteness
condition: for each open subgroup use its permutation module and
Shapiro to embed continuous homomorphisms to \(\mathbf F_\ell\)
in a finite cohomology group.  The local version uses the same
argument with local finite cohomology.

FP2 therefore constructs universal continuous matrix lifts.  For the
disconnected group use its faithful six-dimensional block embedding:
\(\operatorname{diag}(A,\theta A)\) on the identity component,
and the corresponding two off-diagonal blocks on the other
component.  Impose the polynomial block identities and the given
component character, so the resulting maps represent exactly
\(\mathcal G\)-valued lifts.  Quotient by the killed-inertia
entries, by the image of \(I_t\), and by the eight equations
\[
 r(c)_{00}-1,\ r(c)_{11}-1,\ r(c)_{01},\ r(c)_{10},\
 r(c)_{02},\ r(c)_{20},\ Z_{01},\ Z_{10},
 \quad Z=r(g_1)-r(cg_1^{-1}c).
 \tag{OCH.8}
\]
The marked hypotheses are \(c^2=1\), \(\eta(c)=-1\),
\(\Psi(c)=1\), \(\eta(g_1)=1\), and
\(\Psi(g_1)^2-1\in\mathcal O^\times\).
The base point annihilates these equations.  The matrix relation
\(r(c)\theta(r(c))=1\), together with the first six equations
and their residual values, forces \(r(c)=1\); the two remaining
equations are precisely the source normalization.
Quotienting a complete Noetherian local ring by these closed ideals
preserves the required ring category.  Its universal property holds
for all original complete local targets, with uniqueness of the
actual map, proving the assertion.
% END OCH9

% BEGIN OCH10
\paragraph{OCH10. The exact relative presentation and its five weight counts.}
\label{claim:OCH10}
The global box ring before (OCH.8) has a weighted presentation over
\(U/I_t\), with new variables \(z_\nu\) that are actual
matrix-entry deviations.  If \(a_j\) counts the new variables
of weight \(j\), its relation count is
\[
        b_j=a_j+\begin{cases}1&|j|\le1,\\0&|j|>1.\end{cases}
 \tag{OCH.9}
\]
All variables and relations have weights between \(-2\) and \(2\).

\emph{Proof.}
Use the finite set from OCH9 and let
\(N_q=I_q\cap\ker\bar r_0\) at each place away from
\(\ell\).  It is closed and normal in the decomposition group.
Its index in inertia divides the order of the residual image,
which divides \(2|k^\times|\), and hence is prime to \(\ell\).
The local deformation condition of killing \(N_q\) is smooth
in the framed sense used by FP8.  To see the lifting statement,
average over its finite inertia image to conjugate a lifted action
back to its specified finite action.  Invariant matrices are the
image of the same averaging projector; lifts of an invariant vector
can be averaged.  The remaining Frobenius is an intertwiner of two
such finite actions.  Its corrections lie in their invariant
endomorphisms, whose units lift across nilpotent quotients.  This
proves both the small-extension lifting and the corresponding
unramified tangent condition.

Apply FP8's literal presentation theorem and DX3's framed Euler
identity weight by weight.  The possible extra global-dual-invariant
term is zero.  Indeed the \(GD.eta\) condition supplies an element
\(h\) killed by the mod-\(\ell\) cyclotomic character and by
the relevant prime-to-\(\ell\) character component, with
\(\eta(h)=-1\).  Then \(g=ch\) has
\(\eta(g)=1\), \(\bar\Psi(g)=1\), and
\(\bar\epsilon(g)=-1\).  It acts trivially on every residual
adjoint weight and as \(-1\) on its Tate dual.  Since \(\ell\)
is odd there is no invariant.  Complex conjugation acts on matrices
by \(-JX^tJ\).  Its fixed subspace in weight \(j\) has
dimension one for \(j=-1,0,1\), and zero for \(j=-2,2\):
in weights \(\pm1\) it pairs the two entries with a minus
sign; in weight zero it fixes \(\operatorname{diag}(1,0,-1)\);
in weights \(\pm2\) it is minus the identity.  This gives
(OCH.9).  FP8 retains the original local variables and relations,
and chooses each new variable among the actual framed entries;
there is no abstract change to a merely isomorphic deformation
problem.  The six-dimensional faithful embedding has the same
entry-weight differences, so no other weight occurs.
% END OCH10

% BEGIN OCH11
\paragraph{OCH11. Construction of the actual graded model and its numerical shape.}
\label{claim:OCH11}
Every normalized ring \(S\) in OCH9 has the actual graded
presentation, base point and numerical shape required by the original
\texttt{Model} definition.  Its three extension fields
\texttt{pts}, \texttt{der}, and \texttt{lift} are completed
in OCH12--OCH14.  The numerical parameters may be chosen as
\[
 r=a_{-1}+2,\quad s=a_{-2}+1,\quad
 r'=a_1+3,\quad s'=a_2+1,\quad m=a_0+3.
 \tag{OCH.10}
\]
Its ambient smooth chart has weight dimensions
\((s,r,m,r'+1,s'+1)\); its equations have counts
\((s,r+1,m,r',s')\), in weight order \(-2,-1,0,1,2\).
In particular \(r\ge1\).

\emph{Proof.}
Take the ambient complete ring
\(R=(U/I_t)[[z_\nu]]\) before imposing the global relations.
Here is its integral weighted presentation.  Choose finitely many
actual local entry deviations generating the cotangent space of
\(U\) over \(\mathcal O\); successive approximation and
Nakayama give a continuous surjection
\[
 P=\mathcal O[[X_1,\ldots,X_d,z_\nu]]
       \longrightarrow (U/I_t)[[z_\nu]]=R.
\]
Give each \(X_i\) its matrix-entry weight and the new
\(z_\nu\) the weights in OCH10.  If \(J\) is its kernel,
then \(J\) is homogeneous in the precise formal-series sense:
every projection \(f_j\) of an \(f\in J\) onto the
monomials of weight \(j\) also belongs to \(J\).
To prove this integrally, compose the local part with OCH7's
injective map \(U/I_t\hookrightarrow E[[Y_1,\ldots,Y_{12}]]\),
and extend it by the free \(z\)-variables.  This extension is
injective coefficient by coefficient.  Each substituted entry is
homogeneous of its stated weight and has zero constant coefficient,
so substitution commutes with weight projection at every finite
total degree.  If the full image is zero, the image of every
\(f_j\) is zero.  The kernel is precisely \(J\), which
proves the assertion without dividing by any difference of integral
torus eigenvalues.  Completeness and separatedness justify passage
from finite-degree truncations to full series.

Let \(R_j\) be the image of the weight-\(j\) series in
\(P\).  The homogeneous-kernel assertion makes these pieces
independent in finite sums, and multiplication sends
\(R_iR_j\) into \(R_{i+j}\).  Thus
\(A=\bigoplus_j R_j\) is the actual subalgebra of \(R\)
of finite sums of homogeneous series.  Let
\(q:A\to S\) be the induced map.  The relation list \(Q\)
consists of OCH10's relations and (OCH.8), lifted homogeneously to
\(A\).  The eight added weights are
\(0,0,1,-1,2,-2,1,-1\).  Combining them with (OCH.7) and
(OCH.9) gives exactly (OCH.10) and the two asserted count lists.
The base map kills the zero-weight equations as well.

Homogeneous projection of an expression in the complete relation
ideal shows
\(\ker q=(Q)\) in \(A\): each relation list is finite,
and projecting its coefficient series to the complementary weight
retains the required homogeneous expression.  The chart of OCH7,
with the new free coordinates added, gives the required
\texttt{SmoothPt}; its cotangent injectivity is precisely OCH7's
scalar-cleared injectivity.  The weight-zero algebra is a complete
Noetherian local \(\mathcal O\)-algebra.  For completeness,
its presentation uses the finitely many minimal weight-zero
monomials; every weight-zero monomial is a product of these,
by the componentwise well-ordering of \(\mathbf N^d\).
The same argument in each fixed weight shows that \(A\) is
generated over \(A_0\) by finitely many monomials of weights
\(-2,-1,1,2\).  All nonzero weights vanish at the base.
These are the exact \texttt{IsWRing} fields.

Each universal three-dimensional entry has its homogeneous
preimage in \(A\).  The extra six-dimensional entries create
no new condition: after \(r(c)=1\), inverse-transpose entries
are entries at the conjugated group element.  Truncating a
homogeneous formal series modulo the \(n\)-th maximal-ideal
power expresses its image under \(q\) as an
\(\mathcal O\)-linear combination of monomials of actual
entries of the same weight.  Thus both \texttt{ent} and the
all-\(j,n\) field \texttt{dense} hold.  The remaining three
map properties are proved in OCH12--OCH14, on this same \(A,q\).
% END OCH11

% BEGIN OCH12
\paragraph{OCH12. Finite weight zero and actual local representatives of all graded points.}
\label{claim:OCH12}
If the original \texttt{FinA} holds, then
\(M2\) holds for the model of OCH11.  More generally, whenever
\(M2\) holds, every algebraic-closure point of \(A/(Q)\)
can be moved by \(T(t)\) to the restriction of an actual
strictly local point of \(S\).

\emph{Proof.}
The image of \(A_0\) in \(S\) is in the closure of the
weight-zero entry algebra, by OCH11's density.  Under
\texttt{FinA} that algebra is a finite \(\mathcal O\)-module.
It is closed: it is the image of a finite power of the compact ring
\(\mathcal O\) in the Hausdorff complete local ring \(S\).
Thus the image of \(A_0\), a submodule of a finite module, is
finite.  The equality \(\ker q=(Q)\) identifies this image
with the weight-zero ring \((A/(Q))_0\); it is this quotient
ring that is finite over \(\mathcal O\), as required by
\(M2\).  The ambient \(A_0\) need not be finite over
\(\mathcal O\).

Assume now only \(M2\).  Each fixed-weight piece of \(A/(Q)\)
is finite over its weight-zero piece: choose finitely many minimal
monomials of that weight under divisibility by weight-zero monomials.
The quotient weight-zero ring is a finite local
\(\mathcal O\)-algebra.  Every \(\mathcal O\)-linear map
from it to a finite coefficient field is bounded, by choosing finitely
many module generators; hence it is continuous for its adic topology.
The same observation applies to each finite fixed-weight module and
to an \(\mathcal O\)-linear point derivation.  The adic and
uniformizer topologies agree on this finite local quotient, since
its maximal ideal is nilpotent modulo the uniformizer.
Under every algebra map to
\(\overline{\mathbf Q}_\ell\), its elements are integral and
its maximal-ideal elements have absolute value less than one,
since a power of their ideal is contained in the uniformizer ideal.
Choose the finite homogeneous generators of nonzero weights.
For a positive-weight generator of weight \(a\) and a
negative-weight generator of weight \(-b\), the balanced
monomial of powers \(b,a\) is in this maximal ideal.  Whenever
both images are nonzero this gives
\[
              b\,v(\xi x)+a\,v(\xi y)>0.
\]
Therefore the finitely many strict inequalities
\(v(\xi x)+\mathrm{wt}(x)v(t)>0\) admit a common solution;
for generators of only one sign there is only one side to bound.
Choose a rational solution, hence \(t\) in the algebraic closure.
Zero images impose no constraint.  All transformed coordinates now
have absolute value less than one.

Formal evaluation converges in a finite coefficient extension:
the finitely many generator images and a basis of the finite
weight-zero image lie in one such extension.  It agrees with
\(t\cdot\xi\) on all homogeneous series, not just polynomial
coordinates.  In a fixed weight, tails of increasing total degree
are divisible by arbitrarily high powers of the weight-zero maximal
ideal, after extraction of one of finitely many minimal monomials.
Their images tend to zero.  The resulting map annihilates every
complete relation and hence factors through the original \(S\).
Every element of its maximal ideal has image of absolute value
less than one.  This is precisely \texttt{Model.pts}.
% END OCH12

% BEGIN OCH13
\paragraph{OCH13. Extension of every point derivation to the actual complete ring.}
\label{claim:OCH13}
Under \(M2\), let \(f:S\to\overline{\mathbf Q}_\ell\)
be strictly local.  Every point derivation of \(A\) at
\(f\circ q\) annihilating \((Q)\) is the second component
of an actual map
\(S\to\overline{\mathbf Q}_\ell[\varepsilon]/(\varepsilon^2)\)
whose first component is exactly \(f\).

\emph{Proof.}
First the literal strict-local hypothesis already makes the supplied
\(f\) continuous.  If \(u\in S\) had \(|f(u)|>1\),
then \(\ell u^h\in\mathfrak m_S\) for every \(h\),
while \(|f(\ell u^h)|\) would eventually exceed one, a
contradiction.  Thus \(f(S)\) is bounded by one.  A finite
set of generators of \(\mathfrak m_S\) has images bounded
by one number \(q_0<1\), so
\(|f(\mathfrak m_S^h)|\le q_0^h\).  This proves continuity
without an extra hypothesis on the supplied algebra homomorphism.
The map \(a\mapsto f(q(a))+\varepsilon D(a)\) is an algebra
map on \(A\).  Its coordinates have strictly small first
components.  For a formal monomial, its second component is the
finite sum obtained by differentiating one factor.  Along a tail
these terms tend to zero: the first components decay geometrically,
and the possible multiplicity factor is an integer of
\(\ell\)-adic absolute value at most one.  The finitely many
coordinate derivatives lie in one finite extension.
Agreement on arbitrary homogeneous series follows by the fixed-weight
finite-module argument in OCH12, applied also to the derivation.
On a finite \(\mathcal O\)-module every \(\mathcal O\)-linear
map to the finite coefficient field is bounded; the Leibniz rule
makes derivatives of the high powers of the weight-zero maximal
ideal tend to zero as well.  Thus formal evaluation gives a map
on the ambient complete presentation.  It kills its local relations
and the complete ideal generated by \(Q\), because it kills
those relations before completion and multiplication and limits
preserve zero.  Its first component is \(f\): the two maps
agree on the dense polynomial coordinates and are continuous.
This proves the literal \texttt{Model.der}.
% END OCH13

% BEGIN OCH14
\paragraph{OCH14. The integral weight-minus-one lifting field of the model.}
\label{claim:OCH14}
Let \(a\in\mathcal O\setminus\{0\}\).  Suppose a point
derivation \(D:S\to\mathcal O/(a)\) at the base has entry
weight \(-1\), and its local restriction is the reduction of an
integral weight-\(-1\) derivation of \(U\) annihilating
\(I_t\).  Then \(D\circ q\) is the reduction of an integral
weight-\(-1\) point derivation of \(A\).

\emph{Proof.}
Work in the actual ambient ring \((U/I_t)[[z_\nu]]\) of
OCH11, before the global equations.  Use the supplied integral
derivation on its local coefficient ring.  For a new variable of
weight \(-1\), lift its prescribed value from \(\mathcal O/(a)\)
to \(\mathcal O\); put zero on each other variable.
The resulting derivation on formal series is simply its constant
and linear coefficient formula at the base, so it is defined
integrally on the complete ring.  It annihilates the old local
relations by the supplied hypothesis, and restricts to the desired
weight on \(A\).
The reduction agrees with \(D\circ q\) on every coordinate.
For a six-dimensional coordinate this follows by differentiating
\(A^{-1}\) and \(JA^{-t}J\): at the diagonal base its derivative
is a unit scalar multiple of the corresponding three-dimensional
entry derivative with the same weight.  Such coordinates are the
actual ones in FP8.  Equality on coordinates implies equality on
the full complete ring.  Indeed both derivations give local maps
to the complete square-zero ring
\(\mathcal O\oplus\mathcal O/(a)\), hence are adically
continuous, and polynomials in the coordinates are dense.
Restriction proves \texttt{Model.lift}.  Together with
OCH11--OCH13 this completes every field of the original model;
no global equation was required to lift integrally, only its
reduction was required to vanish.
% END OCH14

\subsection{Actual points, rigidity, and admissibility}

% BEGIN OCH15
\paragraph{OCH15. Local conditions of the actual strictly local points.}
\label{claim:OCH15}
Every strictly local point of \(S\) gives the full original
\texttt{IsPt}.  Its representation is conjugate to the base on
inertia away from \(\ell\), after restriction to \(H\).
At every place above \(\ell\), its \(H\)-restriction has a
stable ordered flag with inertia characters
\(\Psi,1,\Psi^{-1}\), and is potentially crystalline.
The same local crystalline condition holds for every actual
first-order deformation through \(S\).

\emph{Proof.}
A strictly local map is adically continuous.  Choose finitely many
generators of the maximal ideal.  Their images have a common norm
bound \(q<1\), while the whole ring has norm at most one:
reduce an element modulo the maximal ideal and lift its finite
residue value by an integral coefficient.  Thus the image of its
\(n\)-th power has norm at most \(q^n\).
Evaluation of a finite formal presentation then lands in the
complete finite field generated by its finitely many coordinate
images.  This proves continuity, integrality and the strict
residual inequalities in \texttt{IsPt}.  The other fields follow
from the actual quotient equations in OCH9.

Away from \(\ell\), both the point and the base factor, on
\(I\cap H\), through the same finite prime-to-\(\ell\)
character quotient \(C\).  The averaging matrix
\[
              M=\sum_{g\in C}r_0(g)^{-1}r(g)
\]
intertwines them.  It is \(|C|I\) plus a matrix all of whose
entries have norm less than one, so it is invertible.  This proves
the actual inertia conjugacy, not merely equality of traces.

For the full inertia at \(\ell\), first work in the formal
chart of OCH7 modulo every power of its variable ideal.  OCH5
makes each truncated representation crystalline, and OCH4 makes
it ordinary.  Its diagonal inertia characters equal their fixed
base characters: the quotient character has weight zero and is
crystalline, so its infinitesimal cocycles are unramified; induction
on the nilpotent length proves equality on inertia.  Hence the
characteristic-polynomial and ordered-product equations hold for
\emph{every} inertia element, including those outside \(D.HL\).
Their images vanish at every order of the formal chart, thus vanish
in that chart and belong to \(I_t\).  They therefore hold at
every point of \(U/I_t\).
Choose an inertia element with three distinct \(\Psi^{w_i}\)
values.  Its spectral projectors and the ordered-product equations
produce the flag, exactly as in OCH5.  Any decomposition element
preserves it: conjugating the chosen inertia element gives another
upper triangular matrix with the same distinct ordered diagonal,
so the intertwiner must be upper triangular, by downward solution
of its matrix entries.  OCH5 gives crystallinity on \(D.HL\).
Conjugate decomposition and inertia groups give every other place;
if the conjugating global element has component \(-1\),
\(\theta\) preserves upper triangular matrices and reverses
and inverts their diagonal, leaving \((\Psi,1,\Psi^{-1})\)
in the same order.  Thus no place or flag ordering is lost.
For an arbitrary algebra map to dual numbers with strictly local
first part, its second component is also controlled by finitely many
coordinates; continuity is not an additional hypothesis.  Extend
\(\mathcal O\) finitely to \(\mathcal O'\) containing
the first coordinate values \(a_i\), all strictly small.
The scalar extension of the complete presentation is still complete,
because \(\mathcal O'\) is finite over \(\mathcal O\).
Formal Taylor division in \(\mathcal O'[[X_i]]\) expresses
every series, modulo the square of \((X_i-a_i)_i\), as its
value at \(a\) plus its linear Taylor terms.  These convergent
identities are ring identities.  The supplied map kills that square
and acts by the prescribed embedding on \(\mathcal O'\),
so it is determined by the finitely many values
\(F(X_i)=a_i+\varepsilon b_i\).  All \(a_i,b_i\) lie in
one finite field, and the Taylor formula proves adic continuity.
Apply OCH5 to this finite algebra.  It includes the nilpotent
directions without passing to their semisimplification.
% END OCH15

% BEGIN OCH16
\paragraph{OCH16. Two-sided points and their actual torus orbits.}
\label{claim:OCH16}
At a strictly local point, being neither upper nor lower triangular
on \(H\) forces exactly one of the following source shapes:
\texttt{TypeI}, irreducibility on \(H\); or \texttt{TypeII},
a middle line and an irreducible outer two-dimensional block.
If two such point representations are isomorphic on \(H\),
they are conjugate by \(T(t)\) on all of \(G\).

\emph{Proof.}
Normalization gives
\[
 r(g_1)-r(cg_1^{-1}c)=\zeta\operatorname{diag}(1,0,-1).
 \tag{OCH.11}
\]
Here \(\zeta\) is congruent, in the strict norm sense, to
\(\Psi(g_1)-\Psi(g_1)^{-1}\), a unit.  Thus \(\zeta\ne0\).
For two points \(\zeta'\ne-\zeta\), because otherwise twice
that fixed unit would have norm less than one, impossible when
\(\ell\ne2\).

An invariant subspace on \(H\) is stable under (OCH.11).
Its three distinct eigenvalues force it to be a coordinate subspace:
apply the three Lagrange interpolation projectors.  A zero entry
\((i,j)\) forced by such a subspace forces also the zero entry
\((2-j,2-i)\), since
\(r(cg^{-1}c)=Jr(g)^tJ\) on \(H\).
Enumerate the six nonempty proper coordinate subsets.
The subsets \(\{0\},\{0,1\}\) force upper triangularity;
\(\{2\},\{1,2\}\) force lower triangularity.
The remaining subsets \(\{1\},\{0,2\}\) force all four
odd-weight entries to vanish.  If the outer block is reducible,
its two distinct marked eigenvalues force a coordinate line and
hence again upper or lower triangularity.  This proves the shapes.

An intertwiner between two points intertwines (OCH.11), so the
nonopposition of \(\zeta,\zeta'\) makes it diagonal,
\(\operatorname{diag}(m_0,m_1,m_2)\), with all \(m_i\ne0\).
Applying the intertwining equation also at \(cg^{-1}c\) gives
\[
 (m_im_{2-i}-m_jm_{2-j})r(g)_{ij}=0.
\]
If \(m_0m_2=m_1^2\), take \(t=m_0/m_1\).
Otherwise all four odd-weight entries vanish, and choose
\(t^2=m_0/m_2\) in the algebraic closure.  In either case
\(r'(g)_{ij}=t^{w_i-w_j}r(g)_{ij}\) on \(H\).
Both representations send \(c\) to the same disconnected element
with matrix one, which commutes with \(T(t)\) in \(\mathcal G\).
Since \(G=H\sqcup Hc\), the equality holds on all of \(G\).
% END OCH16

% BEGIN OCH17
\paragraph{OCH17. The only point having both triangular signs is the base point.}
\label{claim:OCH17}
If an original point is both upper and lower on \(H\), it is
exactly \(r_0\) on \(G\).

\emph{Proof.}
Its three diagonal entries on \(H\) are characters \(a_i\),
and \(a_i/\Psi^{w_i}\) has every value in
\(1+\mathfrak m_{\overline{\mathbf Q}_\ell}\).
By OCH15 the inertia characteristic polynomial is the base one.
Away from \(\ell\), reduction is injective on the finite
prime-to-\(\ell\) values of \(\psi\), so this strict
congruence identifies the correct root:
\(a_i=\Psi^{w_i}\) on inertia.  At \(\ell\), \(GD.gen\)
and the original inertia condition give an inertia element \(s\)
with \(\Psi(s)^2-1\) a unit.  Of the exponents \(1,0,-1\),
only \(w_i\) can be strictly congruent to the \(i\)-th root
at \(s\).  The same is true at an element \(g\) for which
\(\Psi(g)^2-1\) is a unit.  Otherwise apply it to \(gs\)
and divide the character identities at \(gs\) and \(s\).
Thus the same equality holds on all inertia above \(\ell\).

Each ratio is now everywhere unramified over \(K\).  Global
class-field theory makes it factor through the finite ideal class
group.  Its order divides \(h_K\), prime to \(\ell\) by
\(GD.l\_cop\), whereas a root of unity of order prime to
\(\ell\) strictly congruent to one must be one (factor
\(X^n-1\) at \(1\), whose derivative \(n\) is a unit).
All ratios are trivial.  Equality on \(H\), together with
normalization at \(c\), gives equality on \(G\).
% END OCH17

% BEGIN OCH18
\paragraph{OCH18. Type-I points remain irreducible on the full cyclotomic kernel.}
\label{claim:OCH18}
If a point is of type I, its restriction to
\(H\cap\ker\epsilon\), the full infinite cyclotomic kernel,
is irreducible.

\emph{Proof.}
We use an elementary Clifford argument, spelling out its infinite
quotient scope.  The restriction of an irreducible finite-dimensional
representation to a normal subgroup is semisimple: the sum of the
conjugates of a minimal nonzero invariant subspace is the whole
space, and the sum of simple submodules is semisimple.  If the
restriction is reducible, its commuting algebra has dimension
larger than one.  The quotient here is abelian, so conjugation
acts by commuting operators on that algebra.  Its nonzero trace-zero
part is stable.  Over the algebraic closure a commuting family has
a common eigenvector (choose a minimal nonzero common invariant
subspace and an eigenspace inside it).  This yields
\(M\rho(g)=\chi(g)\rho(g)M\), with \(\chi\) trivial on
the normal subgroup.  Its character is nontrivial, since otherwise
Schur's lemma makes the trace-zero matrix scalar and hence zero.
The kernel and image of \(M\) are invariant, so \(M\) is
invertible.  This argument needs no finite-index assumption.

For \(g\) in inertia above \(\ell\) with \(\chi(g)\ne1\),
the trace of \(\rho(g)\) vanishes.  OCH15 gives
\(\Psi(g)+1+\Psi(g)^{-1}=0\), hence \(\Psi(g)^3=1\).
Cyclotomic inertia surjects onto \(\mathbf Z_\ell^\times\):
each finite cyclotomic local extension is totally ramified, and
compactness passes from all finite levels to their inverse limit.
If \(\chi\ne1\), choose an inertia element \(s\) in a
nontrivial \(\chi\)-coset and an inertia element \(u\) with
\(\epsilon(u)=1+\ell\).  If \(\chi(u)\ne1\) apply the
preceding cube relation at \(u\); otherwise apply it at \(su\)
and \(s\) and divide.  In either case \(\Psi(u)^3=1\).
Raise this equality to the finite order of \(\psi\).  It says
\((1+\ell)^{3e\operatorname{ord}\psi}=1\), impossible
in characteristic zero with \(e\ge1\).  Thus no such
self-twist exists, and the restriction is irreducible.
% END OCH18

% BEGIN OCH19
\paragraph{OCH19. The nonpositive base adjoint has no global finite cocycles modulo coboundaries.}
\label{claim:OCH19}
For every \(j\le0\), a global Bloch--Kato cocycle for
\(\operatorname{ad}r_0\), valued in matrix weight \(j\),
is a coboundary, with the exact local conditions of \texttt{rigid0}.

\emph{Proof.}
The involution \(X\mapsto JX^tJ\) preserves each matrix
weight.  Project onto its plus and minus parts using \(1/2\).
The resulting scalar characters are \(\eta\Psi^j\) and
\(\Psi^j\), respectively.  These projections intertwine the
full disconnected action and preserve all local primitives by
KTR1.  If \(j<0\), both are a finite character times a
strictly negative cyclotomic power.  Their finite local condition
at \(\ell\) is zero by the negative-weight Bloch--Kato
formula.  Their away conditions already are inertia coboundaries.
RV7's exact \texttt{soBar} therefore makes each projected global
cocycle a coboundary, on the supplied closure and coefficients.

For \(j=0\), restrict to \(H\).  The scalar characters
\(1,\eta\) become trivial, and the finite condition at
\(\ell\) is unramified by BX7.  Class-field theory makes an
everywhere unramified additive cocycle factor through the finite
class group, so it is zero.  The remaining quotient is of order
two; averaging its cocycle gives a coboundary.  Reassemble the two
matrix projections, using the same equivariant inclusions.
Finite coefficient descent and individual coefficient retractions
are legitimate by KTR1; no uniform field containing all local
primitives is assumed.  This proves the original \texttt{rigid0}
for every nonpositive integer, including weights whose matrix
space is zero.
% END OCH19

% BEGIN OCH20
\paragraph{OCH20. All normalized tangent cocycles satisfy the original global finite conditions.}
\label{claim:OCH20}
If \(F:S\to\overline{\mathbf Q}_\ell[\varepsilon]/(\varepsilon^2)\)
has strictly local first part \(f\), then
\[
       \gamma(g)=\operatorname{snd}(F(r(g)))\,
                         f(r(g))^{-1}
 \tag{OCH.12}
\]
is a continuous global Bloch--Kato cocycle for
\(\operatorname{ad}(r_f)\), with the source's normalization.
Moreover, if this cocycle is a coboundary
\(X-\operatorname{ad}(r_f)(g)X\), then
\(X\in\overline{\mathbf Q}_\ell\operatorname{diag}(1,0,-1)\).

\emph{Proof.}
Differentiate the multiplication in \(\mathcal G\), including
its inverse-transpose component.  The product rule gives the
adjoint cocycle equation; (OCH.8) gives its normalization.  Its
continuity follows from OCH15's formal Taylor division for an
arbitrary supplied dual-number map; it does not presuppose continuity
of its second component.  Away from \(\ell\) it vanishes
on the killed inertia kernel.  Average on the remaining finite
inertia quotient to obtain a coboundary there.  At \(\ell\),
OCH15 makes the dual-number representation crystalline on
\(D.HL\); the extension criterion says exactly that (OCH.12)
has a crystalline period primitive.  Average over the finite
quotient and transport the original decomposition-group graph by
KTR1.  This is the full original global local condition.

For the last assertion put \(H_0=\operatorname{diag}(1,0,-1)\).
The derivative at \(c\) gives \(X=-JX^tJ\).
Differentiating the two marked zero entries in (OCH.11) gives
\([X,\zeta H_0]_{01}=[X,\zeta H_0]_{10}=0\).
Since \(2\) and \(\zeta\) are invertible, the first relation
kills \(X_{11},X_{02},X_{20}\) and pairs the remaining opposite
entries with a minus sign; the marked relations kill
\(X_{01},X_{10}\), hence also \(X_{12},X_{21}\).
Finally \(X_{22}=-X_{00}\).  Thus \(X=X_{00}H_0\).
This is an entrywise calculation over the actual coefficient ring
whenever \(2\) and \(\zeta\) are units.
% END OCH20

% BEGIN OCH21
\paragraph{OCH21. The exact point-rigidity hypotheses follow from InpZ and the marked automorphy input.}
\label{claim:OCH21}
Fix one \(D,g_1\) for which the original \texttt{AutB} and
\texttt{LemW} hold.  Then the original \texttt{PtHyp} holds for
every associated \(U,S\): the fixed-point assertion, every
nonpositive-weight base rigidity assertion, and full rigidity at
every two-sided point.

\emph{Proof.}
The first two assertions are OCH17 and OCH19.  At a two-sided
point use OCH16.  In type I, \texttt{AutB} supplies exactly the
automorphic conjugate-self-dual representation
\(\rho_K\epsilon^{-1}\), for the given \(\iota\).
OCH18 supplies its irreducibility on the full cyclotomic kernel;
twisting does not change the restriction there.  Its pairing is
\(J\), since \(r(c)=1\) gives
\(\rho_K(cgc^{-1})=J\rho_K(g)^{-t}J\).
After the twist the pairing multiplier is \(\epsilon^{-2}\),
which is \(\epsilon^{1-3}\), and its extension to \(c\)
on the adjoint is precisely \(-JX^tJ\).
Thus all hypotheses, signs and multipliers of ADK6's literal
\texttt{AdVanishK} match.  Scalar twisting cancels in the
adjoint.  Every cocycle supplied by OCH20 is a coboundary.

In type II, KTR4--KTR8 prove the actual middle character is one
and construct the continuous irreducible outer extension with
its determinant and full cyclotomic condition.  The same marked
\texttt{LemW} therefore applies; its universal quantifier supplies
modularity for \emph{every} original \texttt{IsExt} extension.
KTR13 is exactly \texttt{RigII} on the full polarized
nine-dimensional adjoint, so again each OCH20 cocycle is a
coboundary.  These are the actual predicates in \texttt{InpZ},
constructed in KIZ2.  The sign convention for coboundaries causes
no change: the usual expression \(gX-X\) becomes the source's
\(X-gX\) by replacing \(X\) by \(-X\).
% END OCH21

% BEGIN OCH22
\paragraph{OCH22. Every two-sided point gives an original admissible representation.}
\label{claim:OCH22}
Under the same \texttt{AutB} and \texttt{LemW}, the actual
\(\rho_f:G_K\to\mathrm{GL}_3(\overline{\mathbf Q}_\ell)\)
of each two-sided point satisfies
\[
       \operatorname{Adm3}(\iota,e,N,\rho_f)
       \quad\text{or}\quad
       \operatorname{Adm2}(\iota,e,N|\operatorname{disc}K|,\rho_f).
 \tag{OCH.13}
\]

\emph{Proof.}
In type I, irreducibility is its defining condition, continuity
is OCH15, and the automorphic field is exactly \texttt{AutB}.
At an away place OCH15 diagonalizes inertia with characters
\(\psi^{w_i}\), each killed by fixing the \(N\)-th roots
of unity.  At a place above \(\ell\), its ordered flag has
characters with inertia
\(\epsilon^{e},1,\epsilon^{-e}\) times those same finite
characters, and OCH15 supplies the potentially crystalline
alternative in the original local condition.  This checks every
field of \texttt{Adm3}, at every original prolongation.

In type II, KTR5--KTR7 give the actual outer representation \(W\)
of \(G\), with determinant \(\eta\), and
\(\rho_f\simeq W|_H\oplus1\) by the explicit coordinate
permutation.  \texttt{LemW} gives the original newform of weight
\(2e+1\) attached to \(W\epsilon^{-e}\), with fixed \(\iota\).
Away from \(\ell\), the outer inertia characters are
\(\psi,\psi^{-1}\) when \(\eta\) is unramified.  If
\(\eta\) is ramified, coprimality of the quadratic discriminant
and \(N\) makes \(\psi\) trivial there, and the residual
quotient is of order two; its characters are \(1,\eta\).
Thus all finite characters are killed by fixing the
\(N|\operatorname{disc}K|\)-th roots, exactly the conductor
allowance in \texttt{Adm2}.
At \(\ell\), the eigenline of character \(\Psi\) in
OCH15 has zero middle coordinate, since some inertia element has
\(\Psi^2-1\) a unit.  Its two outer coordinates give a
nonzero eigenline for \(W\).  Normality of inertia makes it
stable under the decomposition group, and the determinant supplies
the second character, of inertia \(\Psi^{-1}\).
Their cyclotomic weights are \(e,-e\).  No ordered ratio is
\(\epsilon\): raise any proposed equality to the finite order
of the finite characters and evaluate an inertia element of
infinite cyclotomic order; it would force \(2e=1\) or
\(-2e=1\).  This supplies the source's nonexceptional-ratio
alternative for \(W\).  All fields of \texttt{Adm2}, including
the actual direct-sum isomorphism, now hold.
% END OCH22

\subsection{The integral Selmer-to-cotangent comparison}

% BEGIN OCH23
\paragraph{OCH23. A canonical normalized representative of each torsion Selmer class.}
\label{claim:OCH23}
Let \(\alpha=\Psi^{-1}\) act on \(E/\mathcal O\), with
its discrete topology, and let \(\mathcal S\) be the original
strict Selmer group \(\texttt{ShaU}\) with family
\(\texttt{shaFam}(\ell N)\).  For every
\(b\in\mathcal O\setminus\{0\}\), its subgroup
\(\mathcal S[b]\) injects additively into the original
\(\texttt{shaCoc}(b)\).

\emph{Proof.}
Since \(\alpha(g_1)-1\) is a unit, replace a cocycle \(z\)
by
\[
 \nu[z](g)=z(g)-(\alpha(g)-1)
                     (\alpha(g_1)-1)^{-1}z(g_1).
 \tag{OCH.14}
\]
Changing \(z\) by a coboundary leaves this expression unchanged.
It is additive in the class, is zero at \(g_1\), and is zero
as a cocycle exactly when the class is zero.  For a class killed
by \(b\), all its values lie in \((E/\mathcal O)[b]\).
Use the exact isomorphism
\(\mathcal O/(b)\simeq(E/\mathcal O)[b]\),
\(\bar a\mapsto a/b+\mathcal O\), and call the resulting
scalar cocycle \(u\).  Its values are locally constant.

The decomposition group at \(\ell\) belongs to the original
\(\texttt{shaFam}(\ell N)\), and every killed inertia element
belongs to one of its inertia groups.  Thus, on each such local
group, \(\nu[z](g)=(\alpha(g)-1)m\) for a local primitive.
On \(\mathrm{Kill}\), \(\Psi=1\), so the cocycle is zero.
At the distinguished local group, choose its original element
\(s\) with \(\Psi(s)^2-1\) a unit.  Then
\((\alpha(s)-1)bm=0\) forces \(bm=0\).  The local primitive
therefore already lies in \((E/\mathcal O)[b]\); it can be
represented in \(\mathcal O/(b)\).  No enlargement of the
annihilator is necessary.

Send \(u(g)\) to
\[
              \gamma(g)=u(g)(E_{10}-E_{21}).
 \tag{OCH.15}
\]
The full adjoint action of the base on this line is \(\alpha\):
\(-JX^tJ\) fixes \(E_{10}-E_{21}\), and diagonal
conjugation multiplies it by \(\Psi^{-1}\).
Thus (OCH.15) is a weight-\(-1\) matrix cocycle.
At \(c\), \(2u(c)=u(c^2)=0\), so \(u(c)=0\).
Together with \(u(g_1)=0\), the cocycle identity also gives
\(u(cg_1^{-1}c)=0\).  Hence all normalization equations hold.
A local scalar primitive \(m\) gives the source-sign matrix
primitive \(X=-m(E_{10}-E_{21})\), since
\(\operatorname{cob}X=X-\operatorname{ad}(r_0)X\).
This verifies every field of \texttt{shaCoc}; its \((1,0)\)
entry recovers \(u\), so the map is injective.  This proves the
literal \texttt{sha\_inj}, with the allowed choice \(a=b\).
% END OCH23

% BEGIN OCH24
\paragraph{OCH24. Matrix cocycles produce actual normalized square-zero deformations.}
\label{claim:OCH24}
For any \(\gamma\in\texttt{shaCoc}(a)\), there is a point
derivation \(D_\gamma:S\to\mathcal O/(a)\) at the actual
base point satisfying
\[
 D_\gamma(r(g)_{ij})=\Psi(g)^{w_j}\gamma(g)_{ij}.
 \tag{OCH.16}
\]
The map \(\gamma\mapsto D_\gamma\circ q\) is additive and
injective, annihilates every equation in \(Q\), and each of
its values lifts to an integral weight-\(-1\) point derivation
of \(A\).

\emph{Proof.}
In the actual complete ring
\(\mathcal O\oplus\mathcal O/(a)\), with square-zero second
summand, define the matrix of the lift by
\((1+\varepsilon\gamma(g))r_0(g)\), and keep component
\(\eta(g)\).  The cocycle identity is exactly the product
identity in \(\mathcal G\); the inverse is obtained by the
same square-zero formula.  Local constancy of \(\gamma\) and
adic continuity of \(\Psi\) give a continuous lift.  The
\texttt{shaCoc} normalization makes (OCH.8) zero, and its killed
inertia condition gives the required killed subgroup.

For the local condition choose the source's weight-\(-1\)
primitive \(X\) over \(\mathcal O/(a)\), and lift its two
entries to a weight-\(-1\) matrix \(Y\) over
\(\mathcal O\).  Conjugate the base representation over
\(\mathcal O\oplus\mathcal O\) by \(1+\varepsilon Y\).
OCH8 guarantees the exact \(I_t\)-condition on this integral
local deformation.  Its reduction has cocycle \(\gamma\)
on the local group.  Hence the global square-zero representation
satisfies every defining condition of \(S\), with an actual
local comparison map.  Universality gives a map from \(S\).
Its first component is the original base point, by uniqueness;
its second component is the desired derivation.  Multiplication
by the base diagonal gives exactly (OCH.16).

The integral local conjugation supplies the integral local derivative
required by OCH14.  Apply OCH14 to obtain the stated integral
lift on \(A\).  Because \(q(Q)=0\), the reduced derivation
annihilates every equation.  Uniqueness on actual entries in
(OCH.16), followed by continuity on the complete presentation,
proves additivity.  If \(D_\gamma\circ q=0\), the homogeneous
preimage of each entry supplied by OCH11 and the unit
\(\Psi(g)^{w_j}\) force every \(\gamma(g)_{ij}=0\).
This proves injectivity, including its original additive rather
than artificially strengthened linear scope.
% END OCH24

% BEGIN OCH25
\paragraph{OCH25. The full original lower bound, before any finiteness conclusion.}
\label{claim:OCH25}
For every actual model constructed above,
\[
             V\!\operatorname{len}(\Omega,c,e,\psi)
          \le \operatorname{cotLen}(x,A_{-1},Q.F)=M.\operatorname{len}
                   \quad\hbox{in }\mathbf N\cup\{\infty\}.
 \tag{OCH.17}
\]

\emph{Proof.}
For each nonzero \(b\), combine OCH23 and OCH24.  Their
additive injection from \(\mathcal S[b]\) to point derivations
annihilates the weight-\(-1\) equations \(Q.F\), and every
image lifts integrally with that weight.  Nonzero weights vanish
at the base.  These are precisely all hypotheses of SP1, the
proved generic \texttt{LemmaM'}; it gives
\(\operatorname{length}_{\mathcal O}\mathcal S[b]\le M.\operatorname{len}\).

Every class of the original discrete first cohomology group is
torsion.  A continuous cocycle on the compact Galois group has
finite image in the discrete module \(E/\mathcal O\); the
product of finitely many nonzero denominators annihilates it.
Consequently \(\mathcal S\) is the increasing union of its
\(\varpi^n\)-torsion submodules.  The length of an increasing
union is the supremum of these lengths: any finite strict chain
in the union is witnessed by finitely many elements, all in one
member of the union.  Thus the same bound holds for
\(\operatorname{length}_{\mathcal O}\mathcal S\), allowing
infinity on either side.  By RV1's exact quotient and local-family
identification, this is the original
\(V\!\operatorname{len}\), whose level is \(\ell N\).
This establishes \texttt{vlen\_le} without presupposing its
finiteness or the trace upper bound.
% END OCH25

\subsection{The upper bound and the original coefficient datum}

% BEGIN OCH26
\paragraph{OCH26. All five geometric hypotheses on the same model.}
\label{claim:OCH26}
Assume the original \texttt{FinA} and \texttt{PtHyp} for the
same marked datum.  Let \(\sigma\in H\), and suppose that
the trace difference from the base is nonzero at every two-sided
point.  Then the five hypotheses \(M1\)--\(M5\) of SP6 hold
for the model \(A,Q\) and the original trace-difference element.

\emph{Proof.}
OCH12 gives \(M2\).  For \(j=-2,-1,0\), a linear
functional on the rational saturated cotangent quotient annihilating
the weight-\(j\) equations defines a weight-\(j\) point
derivation.  OCH13 lifts it to an actual first-order map on
\(S\).  Its tangent cocycle has that weight because each entry
has a preimage of its prescribed weight.  Base rigidity makes it
a coboundary; OCH20 makes its primitive a multiple of \(H_0\).
But \(H_0\) commutes with the base representation, including
the disconnected action, so this coboundary is zero.  The
derivation is zero on every entry and hence on all of \(A\),
by OCH11's density and the fixed-weight continuity established in
OCH13.  The finite-dimensional cotangent quotient therefore has
zero dual and is zero.  This gives exactly the three saturated
spanning assertions in \(M1\); its smoothness and \(r\ge1\)
were proved in OCH11.

A graded point annihilating both positive and negative weights can
be torus-rescaled to a local point by OCH12.  Its representation
is both upper and lower, hence the base by \texttt{PtHyp}.
Entry density then identifies its graded point with the base;
rescaling does not change it.  This is stronger than the required
\(M3\), so in particular proves that implication with its
original two one-sided alternatives.

For a two-sided graded point, rescale it to a local point and
rescale any two given derivations by the same weight factors.
OCH13 lifts both.  Full point rigidity and OCH20 put both tangent
cocycles in the one-dimensional space of coboundaries of \(H_0\).
The map from these derivations to cocycles is injective by entry
density, so the two derivations are linearly dependent.  Undoing
the common rescaling proves exactly \(M4\).

Finally choose homogeneous diagonal-entry preimages and let
\(a\in A_0\) be their sum at \(\sigma\), minus the base
trace.  Its base value is zero.  Trace is invariant under torus
conjugation; OCH12 therefore transfers the assumed nonzero trace
difference to every two-sided graded point.  This is \(M5\).
All five assertions concern the same \(A,Q,q,a\).
% END OCH26

% BEGIN OCH27
\paragraph{OCH27. A finite irredundant set of actual point orbits and its valuation bound.}
\label{claim:OCH27}
Under OCH26 there is a finite set \(T\) of strictly local,
two-sided maps \(S\to\overline{\mathbf Q}_\ell\), covering
all such point representations up to \(T(t)\)-conjugation,
with no two distinct members so conjugate.  Its model length is
finite and
\[
 \frac{M.\operatorname{len}}{e_{\mathcal O}}
       \le2\sum_{f\in T}v_\ell
        \bigl(\operatorname{tr}\rho_f(\sigma)
                    -\operatorname{tr}\rho_0(\sigma)\bigr).
 \tag{OCH.18}
\]
Here \(e_{\mathcal O}=v_{\mathcal O}(\ell)\), and
\(v_\ell(\ell)=1\).

\emph{Proof.}
Apply SP6, the full \texttt{LemmaM}, with normalized discrete
valuation \(e_{\mathcal O}v_\ell\).  It gives finite model
length and a finite covering set of two-sided graded points.
OCH12 moves each of these to an actual local point.  Both signs
are preserved: vanishing of all negative (respectively positive)
homogeneous coordinates is equivalent to upper (respectively lower)
triangularity of the actual entries.  The forward implication uses
the entry preimages, and the reverse uses their fixed-weight density.
Thus the points obtained are two-sided.  They cover the actual
points as well, because each actual point restricts along \(q\)
to a graded point and equality there determines all actual entries.

On this finite set choose one representative of each actual torus
orbit.  To justify its bound, rather than deleting terms from a
previous sum, restrict this new set again along \(q\).
It is still a full graded cover: lift any graded point by OCH12
and use the actual cover.  Conversely actual torus conjugacy
implies the same graded torus action on every homogeneous element,
by entry density and continuity.  The map \(f\mapsto f\circ q\)
is injective on our representatives, since equality gives equal
entry representations and hence the same orbit.  Apply SP6's
\emph{every-cover} inequality to this restricted set.  Its trace
values are the original ones and its indexing has no multiplicity.
Divide by the positive integer \(e_{\mathcal O}\) to obtain
(OCH.18).  This proves the literal irredundant \texttt{model\_upper}
consumer with all its point and valuation quantifiers.
% END OCH27

% BEGIN OCH28
\paragraph{OCH28. The unramified-coefficient OneChar conclusion.}
\label{claim:OCH28}
Assume \(e_{\mathcal O}=1\), the exact \(GD\) datum, the
given \(\iota\), \(\sigma\), and base trace \(t\), and
assume that every original \(\operatorname{Adm3}\) or
\(\operatorname{Adm2}\) representation has trace at
\(\sigma\) different from \(t\).  Under \texttt{InpZ} and
\texttt{AutomorphyUnramifiedCoefficients}, \(V\!\operatorname{len}\)
is finite and there is a finite set \(X\) of actual admissible
representations such that distinct members are not isomorphic and
\[
             V_n(\Omega,c,e,\psi)
                \le2\sum_{\rho\in X}
                       v_\ell(\operatorname{tr}\rho(\sigma)-t).
 \tag{OCH.19}
\]

\emph{Proof.}
Choose \(D\) by OCH9.  Invoke the original automorphy input
\emph{once}, for this \(D\) and this \(\iota\), obtaining
one marked \(g_1\) with \texttt{FinA}, \texttt{AutB}, and
\texttt{LemW}.  Construct \(U,S\) by OCH9 and their model
by OCH10--OCH14, using the actual \texttt{LemmaP} in
\texttt{InpZ}.  OCH21 supplies \texttt{PtHyp}; OCH22 makes
every two-sided point admissible.  The hypothesis on traces
therefore supplies OCH26's nonvanishing condition.
OCH27 gives the finite irredundant point set \(T\), and OCH25
gives the lower bound.  Its finite upper bound makes the original
Selmer length finite before taking its natural-number value.

Set \(X=\{\rho_f:f\in T\}\), where
\(\rho_f(\gamma)=r_f(\operatorname{galEmb}\gamma).\mathrm{left}\).
These are the actual source homomorphisms.  If two are isomorphic,
OCH16 gives a torus conjugacy of the two full point representations,
and the irredundancy of \(T\) gives equality of the points.
Thus the map \(T\to X\) is injective and \(X\) has the
required isomorphism-irredundancy.  Since \(\eta(\sigma)=1\),
its base trace is precisely the supplied
\(t=\Psi(\sigma)+1+\Psi(\sigma)^{-1}\).
Combine (OCH.17) and (OCH.18).  The source normalization of
\(V_n\) is length divided by \(e_{\mathcal O}\); the sum
is unchanged under the injection \(T\to X\).  This is
(OCH.19), with all three conclusions of the original predicate.
% END OCH28

% BEGIN OCH29
\paragraph{OCH29. An unramified coefficient datum with exactly the same character values.}
\label{claim:OCH29}
For every original character with \(\ell\nmid\operatorname{ord}\psi\),
there is an original unramified coefficient datum \(c'\) and a
Dirichlet character \(\psi'\) of the same modulus having
exactly the same values in the fixed
\(\overline{\mathbf Q}_\ell\).  The full \(GD\) predicate
is preserved.

\emph{Proof.}
Put \(m=\operatorname{ord}\psi\).  Choose in the fixed
algebraic closure a primitive \(m\)-th root of unity.  It lies
in the unramified extension whose residue field contains
\(\mu_m\): lift the simple roots of \(X^m-1\), using
\(\ell\nmid m\).  Its valuation ring is finite unramified
over \(\mathbf Z_\ell\), and hence is an allowed \(c'\)
with ramification index one.  All values of \(\psi\) on units
are powers of that root, and its nonunit values are zero.
Regard these same values in \(c'\); injectivity into the
algebraic closure verifies multiplicativity and every
Dirichlet-character condition.

Order, parity, and all primary components are unchanged.  For the
last statement, primary decomposition is the unique decomposition
of a finite-order element into its prime-primary and prime-to-that-
prime parts; injective coefficient maps preserve both powers and
orders and hence preserve those components.  The Galois characters
are formed by composition with the same cyclotomic restriction
map, so their kernels, including the \(GD.eta\) kernel
condition, are identical after inclusion in the fixed closure.
Consequently \(\operatorname{comp}_\ell(\psi')^2=1\),
\(\Psi'(-1)=\Psi(-1)\), and every value-dependent field of
\(GD\) is unchanged.  All other fields concern only
\(K,\ell,N,e\).  This proves the exact coefficient-transport
assertion, without asserting an embedding of the entire possibly
ramified original \(\mathcal O\) into \(\mathcal O'\).
% END OCH29

% BEGIN OCH30
\paragraph{OCH30. The full original OneChar, on every supplied coefficient datum.}
\label{claim:OCH30}
The literal original implication is proved:
\[
       \texttt{InpZ}\ \wedge\
       \texttt{AutomorphyUnramifiedCoefficients}
                    \quad\Longrightarrow\quad\texttt{OneChar}.
 \tag{OCH.20}
\]
It retains every original closure, coefficient datum, fixed
\(\iota\), \(N\ne0\), \(e\), \(GD\),
\(\sigma\) and \(t\) quantifier.  In particular it adds
no lower bound on \(\ell\) or on \(e\) beyond the original
hypotheses.

\emph{Proof.}
Use OCH29 to pass from \((c,\psi)\) to
\((c',\psi')\) with exactly equal algebraic-closure values.
The two \(\Psi\)'s and their inverses agree at every element;
thus the specified \(t\) is unchanged.  The classes of
\(\operatorname{Adm3}(\iota,e,N)\) and
\(\operatorname{Adm2}(\iota,e,N|\operatorname{disc}K|)\)
representations do not refer to this auxiliary coefficient ring.
The trace-nonvanishing premise therefore stays exactly the same,
with the same \(\iota\) and \(\sigma\).
Apply OCH28 to \(c',\psi'\).  RV7's full coefficient
independence gives simultaneous finiteness of the two original
\(V\!\operatorname{len}\)'s and equality of their normalized
\(V_n\)'s.  Carry the same finite set of actual representations
\(X\) back without changing a matrix or an isomorphism condition.
Its inequality is (OCH.19) for the original coefficient datum.
These are exactly the conclusion fields in
\texttt{KatoUse.lean, OneChar}; the proof never replaces its
irredundant finite set by a formal multiset or assumes finiteness
in order to obtain it.

KIZ2 supplies the \texttt{InpZ} argument.  The separate assertion
\texttt{AUT\_UNRAM\_COEF} remains a premise of (OCH.20).
Accordingly this chapter proposes closure of the selected
\texttt{TASK\_ONE\_CHAR\_PROOF}, whose recorded contract is
this implication.  It does not assert that the separate automorphy
producer, the complete arithmetic installation, or RH has thereby
been completed.
% END OCH30
